\SecondArticle{6}
Consider the product of $m(m-1)$ factors
\[
\begin{aligned}
&(a-b)(a-c)(a-d)\, .\, .\, .\, .\, .\\
\times{}&(b-a)(b-c)(b-d)\, .\, .\, .\, .\, .\\
\times{}&(c-a)(c-b)(c-d)\, .\, .\, .\, .\, .\\
\times{}&(d-a)(d-b)(d-c)\, .\, .\, .\, .\, .\\
&\text{etc.}
\end{aligned}
\]
which we shall denote by $\pi$. Since it involves the indeterminates $a,b,c$ etc. symmetrically, we shall suppose it reduced to the form of a polynomial in $\lambda',\lambda'',\lambda'''$ etc. Let this polynomial become $p$ when $l',l'',l'''$ etc. are substituted respectively for $\lambda',\lambda'',\lambda'''$ etc. With this established, we shall call $p$ the \textit{determinant} of the polynomial
\[
y=x^m-l'x^{m-1}+l''x^{m-2}-l'''x^{m-3}+\text{etc.}
\]
\SourceLeaf{046}
Thus, for example, for $m=2$ we have
\[
p=-{l'}^2+4l''
\]
Likewise, for $m=3$ we find
\[
p=-{l'}^2{l''}^2+4{l'}^3l'''+4{l''}^3-18l'l''l'''+27{l'''}^2
\]
The determinant of the polynomial $y$ is therefore a polynomial in the coefficients $l',l'',l'''$ etc. which, under the substitutions $l'=\lambda',l''=\lambda'',l'''=\lambda'''$ etc., becomes the product of all the differences between pairs of the quantities $a,b,c$ etc. In the case $m=1$, that is, when there is only one indeterminate $a$ and hence no differences at all, it will be convenient to take the number 1 itself as the determinant of the polynomial $y$.

\Needspace{8\baselineskip}In establishing the notion of the determinant, we had to regard the coefficients of the polynomial $y$ as indeterminate quantities. The determinant of a polynomial with fixed coefficients,
\[
Y=x^m-L'x^{m-1}+L''x^{m-2}-L'''x^{m-3}+\text{etc.}
\]
will be a fixed number $P$, namely the value of the polynomial $p$ for $l'=L',l''=L''$, $l'''=L'''$ etc. Thus, if we suppose that $Y$ can be resolved into linear factors,
\[
Y=(x-A)(x-B)(x-C)\, .\, .\, .\, .\, .
\]
or that $Y$ arises from
\[
\upsilon=(x-a)(x-b)(x-c)\, .\, .\, .\, .\, .
\]
by setting $a=A,b=B,c=C$ etc., so that the same substitutions turn $\lambda',\lambda'',\lambda'''$ etc. respectively into $L',L'',L'''$ etc., then $P$ will clearly equal the product of the factors
\[
\begin{aligned}
&(A-B)(A-C)(A-D)\, .\, .\, .\, .\, .\\
\times{}&(B-A)(B-C)(B-D)\, .\, .\, .\, .\, .\\
\times{}&(C-A)(C-B)(C-D)\, .\, .\, .\, .\, .\\
\times{}&(D-A)(D-B)(D-C)\, .\, .\, .\, .\, .\\
&\text{etc.}
\end{aligned}
\]
It follows that, if $P=0$, at least two of the quantities $A,B,C$ etc. must be equal; conversely, if it is not the case that $P=0$, all the quantities $A,B,C$ etc. must be unequal. We now observe that, on setting $\frac{\mathrm dY}{\mathrm dx}=Y'$, or

\SourceLeaf{047}
\[
\begin{gathered}Y'=mx^{m-1}-(m-1)L'x^{m-2}+(m-2)L''x^{m-3}\\-(m-3)L'''x^{m-4}+\text{etc.}\end{gathered}
\]
we have
\[
\begin{aligned}
Y'={}&(x-B)(x-C)(x-D)\, .\, .\, .\, .\, .\\
&+(x-A)(x-C)(x-D)\, .\, .\, .\, .\, .\\
&+(x-A)(x-B)(x-D)\, .\, .\, .\, .\, .\\
&+(x-A)(x-B)(x-C)\, .\, .\, .\, .\, .\\
&+\text{etc.}
\end{aligned}
\]
\Needspace{16\baselineskip}Thus, if two of the quantities $A,B,C$ etc. are equal, for example $A=B$, then $Y'$ will be divisible by $x-A$; in other words, $Y$ and $Y'$ will have the common divisor $x-A$. Conversely, if $Y'$ is supposed to have any common divisor with $Y$, then $Y'$ must contain one of the linear factors $x-A,x-B,x-C$ etc., for example the first, $x-A$. This clearly cannot happen unless $A$ equals one of the remaining quantities $B,C,D$ etc. From all this we therefore obtain two \textsc{Theorems}:

I. \textit{If the determinant of the polynomial $Y$ is $=0$, then $Y$ certainly has a common divisor with $Y'$; hence, if $Y$ and $Y'$ have no common divisor, the determinant of the polynomial $Y$ cannot be $=0$.}

II. \textit{If the determinant of the polynomial $Y$ is not $=0$, then $Y$ and $Y'$ certainly cannot have a common divisor; or, if $Y$ and $Y'$ have a common divisor, the determinant of the polynomial $Y$ must necessarily be $=0$.}

\SecondArticle{7}
But it must be carefully noted that the whole force of this very simple proof rests on the supposition that the polynomial $Y$ can be resolved into linear factors. At this very point, where the general possibility of such a resolution is to be proved, that supposition would be nothing but begging the question. Yet not all those who have attempted analytic proofs of the main theorem have guarded against fallacies exactly like this. We have already pointed out the origin of this plausible illusion in the very statement of the inquiry: all of them investigated only the \textit{form} of the roots of equations, whereas they ought to have proved their \textit{existence}, which they rashly assumed. But enough has already been said in the paper cited above about this way of proceeding, which departs too far from rigour and clarity. We shall therefore now place the theorems of the preceding article,
\SourceLeaf{048}
at least one of which we cannot do without for our purpose, on a firmer foundation. We shall begin with the second, as the easier one.

\Needspace{18\baselineskip}\SecondArticle{8}
Denote by $\rho$ the function
\[
\begin{aligned}
&\frac{\pi(x-b)(x-c)(x-d)\, .\, .\, .\, .\, .}{(a-b)^2(a-c)^2(a-d)^2\, .\, .\, .\, .\, .}\\
&+\frac{\pi(x-a)(x-c)(x-d)\, .\, .\, .\, .\, .}{(b-a)^2(b-c)^2(b-d)^2\, .\, .\, .\, .\, .}\\
&+\frac{\pi(x-a)(x-b)(x-d)\, .\, .\, .\, .\, .}{(c-a)^2(c-b)^2(c-d)^2\, .\, .\, .\, .\, .}\\
&+\frac{\pi(x-a)(x-b)(x-c)\, .\, .\, .\, .\, .}{(d-a)^2(d-b)^2(d-c)^2\, .\, .\, .\, .\, .}\\
&+\text{etc.}
\end{aligned}
\]
which, since $\pi$ is divisible by each denominator, is a polynomial in the indeterminates $x,a,b,c$ etc. Further, set $\frac{\mathrm d\upsilon}{\mathrm dx}=\upsilon'$, so that we have
\[
\begin{aligned}
\upsilon'={}&(x-b)(x-c)(x-d)\, .\, .\, .\, .\, .\\
&+(x-a)(x-c)(x-d)\, .\, .\, .\, .\, .\\
&+(x-a)(x-b)(x-d)\, .\, .\, .\, .\, .\\
&+(x-a)(x-b)(x-c)\, .\, .\, .\, .\, .\\
&+\text{etc.}
\end{aligned}
\]
For $x=a$ we clearly have $\rho\upsilon'=\pi$. Hence we conclude that the polynomial $\pi-\rho\upsilon'$ is identically divisible by $x-a$, and likewise by $x-b,x-c$ etc., as well as by their product $\upsilon$. Thus, setting
\[
\frac{\pi-\rho\upsilon'}{\upsilon}=\sigma
\]
we obtain for $\sigma$ a polynomial in the indeterminates $x,a,b,c$ etc. which, like $\rho$, is symmetric with respect to the indeterminates $a,b,c$ etc. We can therefore find two polynomials $r,s$ in the indeterminates $x,l',l'',l'''$ etc. which, under the substitutions $l'=\lambda',l''=\lambda'',l'''=\lambda'''$ etc., become $\rho,\sigma$ respectively. Following the analogy, let us denote the polynomial
\[
\begin{gathered}mx^{m-1}-(m-1)l'x^{m-2}+(m-2)l''x^{m-3}\\-(m-3)l'''x^{m-4}+\text{etc.}\end{gathered}
\]
that is, the derivative $\frac{\mathrm dy}{\mathrm dx}$, by $y'$, so that $y'$ becomes $\upsilon'$ under those same
\SourceLeaf{049}
substitutions. Then $p-sy-ry'$ clearly becomes $\pi-\sigma\upsilon-\rho\upsilon'$ under the same substitutions, that is, 0, and must therefore already vanish identically in itself (Article 5). We thus have the identity
\[
p=sy+ry'
\]
Hence, if we suppose that substituting $l'=L',l''=L'',l'''=L'''$ etc. gives $r=R,s=S$, we also have identically
\[
P=SY+RY'
\]
Here $S,R$ are polynomials in $x$, while $P$ is a fixed quantity or number. It follows at once that $Y$ and $Y'$ cannot have a common divisor unless $P=0$. This is precisely the second theorem of Article 6.

